\documentclass[12pt,a4paper]{article}

\usepackage[a4paper,margin=1in,top=1.25in]{geometry}
\usepackage{setspace}
\usepackage{lmodern}
\usepackage{microtype}
\usepackage{amsmath,amssymb}
\usepackage{graphicx}
\usepackage{booktabs}
\usepackage{caption}
\usepackage{subcaption}
\usepackage{float}
\usepackage{xcolor}
\usepackage[colorlinks=true,linkcolor=blue,urlcolor=blue]{hyperref}

\graphicspath{{results/curves/}{upload/}}
\setstretch{1.08}

\title{
    \vspace{2cm}
    \textbf{\Huge Systems Engineering for Deep Learning}\\[0.5cm]
    \textbf{\Large Assignment 2}\\[0.2cm]
    \large Course Code: CS6886
}
\author{
    \textbf{\Large Aditya Girish}\\[0.2cm]
    Roll No: \texttt{EP23B048}\\[0.2cm]
    B.Tech Engineering Physics\\[0.2cm]
    IIT Madras
}
\date{September 2026}

\begin{document}

\begin{figure}[t]
    \centering
    \includegraphics[width=0.25\linewidth]{image.png}
\end{figure}

\maketitle
\thispagestyle{empty}
\newpage

\section{Question 1: Training Baseline}

\subsection{CIFAR-10 pipeline}

The CIFAR-10 training set contains 50,000 images and the test set contains
10,000 images, each with spatial resolution $32\times32$ and one of ten
classes. The training pipeline used the following augmentation and
normalization sequence:
\begin{enumerate}
    \item \texttt{RandomCrop(32, padding=4)},
    \item \texttt{RandomHorizontalFlip()},
    \item conversion to a tensor, and
    \item channel-wise normalization using CIFAR-10 mean
    $(0.4914, 0.4822, 0.4465)$ and standard deviation
    $(0.2470, 0.2435, 0.2616)$.
\end{enumerate}
For testing, no random augmentation was applied; images were converted to
tensors and normalized using the same mean and standard deviation. A fixed
seed of 6886 was used for reproducibility.

\subsection{Model and training strategy}

The baseline is torchvision's MobileNetV2 initialized using the public
ImageNet-1K V2 pretrained weights. The first convolutional layer was adapted
to stride 1 (rather than ImageNet's initial stride 2) to preserve useful
spatial resolution for $32\times32$ CIFAR-10 inputs. The original ImageNet
classifier was replaced by a newly initialized linear layer with 10 outputs.
The width multiplier was 1.0, dropout probability was 0.2, and the standard
torchvision BatchNorm2d settings were retained (including momentum 0.1).

The new classifier was trained for the first two epochs while the pretrained
feature extractor was frozen. The entire network was then unfrozen and
fine-tuned end-to-end for the remaining 18 epochs. Training used minibatches
of 128, SGD with Nesterov momentum (momentum 0.9), initial learning rate 0.01,
weight decay $4\times10^{-5}$, and a cosine-annealing learning-rate schedule
over 20 epochs. Cross-entropy loss was used without label smoothing.

\subsection{Results}

Table~\ref{tab:baseline-results} reports the main baseline result. The checkpoint
was selected using the fixed 5,000-image validation split; the official test set
was evaluated only after selection. The supplied checkpoint's held-out test log
records 92.39\% top-1 accuracy.

\begin{table}[H]
    \centering
    \caption{Uncompressed pretrained-MobileNetV2 baseline on CIFAR-10.}
    \label{tab:baseline-results}
    \begin{tabular}{lc}
        \toprule
        Metric & Value \\
        \midrule
        Number of training epochs & 20 \\
        Held-out test top-1 accuracy & \textbf{92.39\%} \\
        Best validation accuracy & 92.86\% \\
        Final-epoch train accuracy & 97.27\% \\
        Final-epoch validation accuracy & 92.78\% \\
        Held-out test loss & 0.2501 \\
        \bottomrule
    \end{tabular}
\end{table}

Figure~\ref{fig:baseline-curves} shows the training and test curves. The
accuracy rises quickly after the pretrained backbone is unfrozen at epoch 3:
validation accuracy increases from 63.20\% at epoch 2 to 85.88\% at epoch 3. It then
improves steadily and plateaus around 92--93\% after approximately 14 epochs.

\begin{figure}[H]
    \centering
    \includegraphics[width=0.94\linewidth]{baseline.png}
    \caption{Training and test loss and top-1 accuracy for the uncompressed
    pretrained MobileNetV2 baseline.}
    \label{fig:baseline-curves}
\end{figure}

\subsection{Failure modes and limitations}

The curves indicate mild overfitting late in training: final training accuracy
is 97.27\% while final validation accuracy is 92.78\%, a gap of approximately
4.5 percentage points. Validation loss is essentially flat late in training,
whereas training loss continues to decrease. Consequently,
substantially extending this schedule is unlikely to improve generalization
without stronger regularization or additional augmentation. A second practical
failure mode is accidental CPU execution in a hosted notebook environment;
training must be run with CUDA available and the recorded device should be
checked before a long run. Finally, this baseline uses ImageNet-pretrained
weights, so its accuracy should be interpreted as transfer-learning performance
rather than accuracy obtained by training MobileNetV2 from random
initialization.

























\end{document}
